\documentclass[11pt]{article}
\usepackage[margin=25mm]{geometry}
\usepackage{booktabs,siunitx,tabularx,array}
\newcolumntype{Y}{>{\raggedright\arraybackslash}X}
\begin{document}
\section*{A table, built in three steps}
All numbers below are invented teaching data for a typesetting lesson.

\subsection*{1. Understand rows and columns}
The first column is left aligned; the second is right aligned.
\begin{center}
\begin{tabular}{lr}
Task & Time (min) \\
Draft & 9.50 \\
Review & 12.25 \\
\end{tabular}
\end{center}

\subsection*{2. Separate the heading and align decimals}
The S column parses numbers. Braces protect its text heading.
\begin{center}
\begin{tabular}{lS[table-format=2.2]}
\toprule
Task & {Time (min)} \\
\midrule
Draft & 9.50 \\
Review & 12.25 \\
\bottomrule
\end{tabular}
\end{center}

\subsection*{3. Add a description without overflowing}
The description receives the remaining width and wraps.
\begin{table}[htbp]
\centering
\caption{Invented tasks and times}\label{tab:tasks}
\begin{tabularx}{\linewidth}{@{}lY S[table-format=2.2]@{}}
\toprule
Task & Description & {Time (min)} \\
\midrule
Draft & Write a short introduction that states the question before explaining the background. & 9.50 \\
Review & Check that captions name units, references resolve, and the table fits its container. & 12.25 \\
\bottomrule
\end{tabularx}
\end{table}
Table~\ref{tab:tasks} keeps the descriptions readable.
Change one description and recompile to see the X column reflow.
\end{document}